% Preámbulo común de los documentos legales (legal/*.tex).
% Se compila cada uno con:   cd legal && latexmk -pdf <archivo>.tex
% La paleta es la del producto y la del manual: si cambia una, cambian las dos.
\usepackage[T1]{fontenc}
\usepackage[utf8]{inputenc}
\usepackage[spanish,es-noshorthands,es-tabla]{babel}
\usepackage[scaled=0.92]{helvet}
\renewcommand{\familydefault}{\sfdefault}
\usepackage{amssymb}
% Los que usan los textos y pdflatex no trae
\DeclareUnicodeCharacter{2713}{\checkmark}
\DeclareUnicodeCharacter{22EF}{\ensuremath{\cdots}}
\DeclareUnicodeCharacter{2192}{\ensuremath{\rightarrow}}
\usepackage[a4paper,top=2cm,bottom=2.2cm,left=2.2cm,right=2.2cm,headheight=14pt]{geometry}
\usepackage[table]{xcolor}
\usepackage{microtype}
\usepackage{enumitem}
\usepackage{booktabs}
\usepackage{tabularx}
\usepackage{array}
\usepackage{longtable}
\usepackage{fancyhdr}
\usepackage{titlesec}
\usepackage{lastpage}
\usepackage[most]{tcolorbox}
\usepackage[hidelinks]{hyperref}

\definecolor{tinta}{HTML}{0B2239}
\definecolor{tinta2}{HTML}{2A4864}
\definecolor{tenue}{HTML}{496A86}
\definecolor{marca}{HTML}{1668C4}
\definecolor{ambar}{HTML}{A67300}
\definecolor{rojo}{HTML}{C2001D}
\definecolor{fondo}{HTML}{F4F7FB}
\definecolor{linea}{HTML}{C3D4E4}
\definecolor{ambarfondo}{HTML}{FDF6E4}
\definecolor{ambarborde}{HTML}{E4C88A}
\definecolor{azulfondo}{HTML}{E7F0FB}
\definecolor{azulborde}{HTML}{C9DDF4}

\color{tinta}
\setlength{\parindent}{0pt}
\setlength{\parskip}{5pt plus 2pt}
\linespread{1.08}
\setlist{itemsep=2pt,topsep=3pt}

\titleformat{\section}{\color{tinta}\large\bfseries}{}{0em}{}
\titleformat{\subsection}{\color{tinta2}\normalsize\bfseries}{}{0em}{}
\titlespacing*{\section}{0pt}{12pt}{4pt}
\titlespacing*{\subsection}{0pt}{8pt}{2pt}

\renewcommand{\arraystretch}{1.25}
\arrayrulecolor{linea}

% Lo que hay que completar a mano: una línea para escribir
\newcommand{\llenar}[1][6cm]{\rule[-2pt]{#1}{0.5pt}}
% Lo que tiene que decidir o completar el abogado o las partes
\newcommand{\completar}[1]{{\color{ambar}[#1]}}
% Una casilla para marcar
\newcommand{\casilla}{\fbox{\rule{0pt}{7pt}\rule{7pt}{0pt}}}

\newtcolorbox{nota}[1][]{colback=azulfondo,colframe=azulborde,boxrule=0.6pt,arc=3pt,
  left=8pt,right=8pt,top=5pt,bottom=5pt,fontupper=\small\color{tinta2},#1}
\newtcolorbox{borrador}{colback=ambarfondo,colframe=ambarborde,boxrule=0.6pt,arc=3pt,
  left=8pt,right=8pt,top=5pt,bottom=5pt,fontupper=\small}

% Pie: qué documento es, su versión y la página
\newcommand{\documento}[2]{%
  \pagestyle{fancy}\fancyhf{}%
  \renewcommand{\headrulewidth}{0pt}%
  \fancyfoot[L]{\scriptsize\color{tenue}#1 · #2}%
  \fancyfoot[R]{\scriptsize\color{tenue}Página \thepage\ de \pageref*{LastPage}}%
}

% Firma: una línea y lo que va debajo
\newcommand{\firma}[1]{\begin{minipage}[t]{0.44\textwidth}\vspace{1.6cm}\rule{\linewidth}{0.5pt}\\{\small\color{tinta2}#1}\end{minipage}}
